\subsection{Nine through fourteen points}

The source files for these six layers are in
\begin{center}
\file{computation/n9_to_14/}.
\end{center}
From that directory run
\begin{verbatim}
python run_verification.py --jobs 8 --seconds 240
\end{verbatim}
The number of workers may be changed without changing the models.  Every
individual solver uses a fixed random seed; a timeout has status
\texttt{UNKNOWN} and is retained.  A successful replay ends with
\file{certificates/manifest.json}, whose seven terminal checks are
\texttt{true} and whose status is \texttt{PASS}.  If a replay is interrupted
after all terminal certificates have been written, the same terminal and
SHA-256 audit can be rerun with \file{audit_manifest.py}.
The active lower-bound table contains $c(7)=11$.  All certificates in the
present package were regenerated after that correction; in particular, the
twelve-point recursion now has 622, rather than 617, nine-point child
count vectors.

The files are grouped by mathematical role as follows.
Some filenames retain implementation terminology: in those names,
\texttt{signature} means the global count vector defined in
Section~\ref{sec:small-uniform-en}, \texttt{profile} means the degrees attached
to point pairs, and \texttt{support} means an actual subset of labelled points.
These words are not additional mathematical notions.
\begin{longtable}{P{6.2cm}P{8.3cm}}
\toprule File & Role in the verification\\ \midrule
\file{signature_filters.py} & enumerates all global count vectors from
\eqref{eq:small-line-pairs-en}--\eqref{eq:small-triple-partition-en}, the
largest-block bounds, deletion averages, and the cited line-arrangement
inequalities\\
\file{projected_arrangement_model.py} & generates all numerical local
multiplicity types and their line/circle splits; the final source contains no
historical line-deletion obstruction table\\
\file{verify_local_type_filter.py} & checks exact unlabelled sums of the local
multiplicity vectors and writes \file{nN_01_signature_filter.json}\\
\file{verify_augmented_line_melchior.py} & checks first
\eqref{eq:augmented-line-summed-en}, then the sums over all points required by
\eqref{eq:augmented-line-point-en}; it writes certificates 02 and 03\\
\file{filter_by_classified_split_endpoints.py} & checks the consistency
identities linking point degrees and point-pair degrees in
\eqref{eq:small-category-totals-en}--\eqref{eq:small-pair-profile2-en};
the runner explicitly disables both optional finite catalogues\\
\file{pair_profile_generator.py},
\file{verify_pair_moment_filter.py} & generate all point-pair degree data and
impose the exact sums of intersection sizes and of their second binomial
coefficients, together with the disjoint-triple counts; they write
certificate 05\\
\file{block_intersection_rows.py},
\file{block_row_existence_filter.py} & at $n=10$, impose the three exact
triple strata around each represented block and the bound on matchings between
two blocks, reducing 21 count vectors to 12\\
\file{verify_recursive_inversion_rows.py} & maps each local multiplicity vector
to the full count vector one order lower, balances all point totals at each
level, and at the nine-point base selects the maximal lines and circles as
subsets of the labelled points\\
\file{verify_n10_final_radical_axis.py} & reconstructs the two matching
orbits, the unique surviving line-design orbit, and the five exact minors that
exclude the final ten-point count vector; it first reads certificate 07, records
its SHA-256 digest, and checks that the retained rows force the two five-point
circles to partition the point set\\
\file{verify_n13_n14_global_inequalities.py} & independently regenerates the
3962 and 8981 raw count vectors and checks the global exclusions and
\eqref{eq:n13-global-two-ineq-en}; it does not select labelled point subsets\\
\file{run_verification.py}, \file{audit_manifest.py} & execute the stages in
order and check all terminal statuses, the unique intermediate $n=10$ count
vector, file sizes, and SHA-256 digests\\
\bottomrule
\end{longtable}

The terminal certificate paths are
{\raggedright
\begin{description}[style=nextline,leftmargin=2.8em]
\item[$n=9$] \file{certificates/n9_04_exact_support.json}.
\item[$n=10$] \file{certificates/n10_07_recursive.json} and
  \file{certificates/n10_08_final_geometry.json}.
\item[$n=11$] \file{certificates/n11_06_recursive.json}.
\item[$n=12$] \file{certificates/n12_06_recursive.json}.
\item[$n=13,14$] \file{certificates/n13_n14_global_inequalities.json}.
\end{description}
\par}
The intermediate filenames have the numerical prefixes displayed in the
table and are consumed in that order.  Each stage after 01 records the path
and SHA-256 digest of its immediate input whenever it reads one.  The final
manifest records forty-two source and certificate files.  Thus a reader may
delete the \file{certificates} directory and regenerate the complete chain
from the listed source files.

In the audited four-worker replay, the full chain took 348.52 seconds of wall
time.  The largest labelled base model had 840 Boolean variables and 546
constraints, and the slowest individual recursive solve took 6.62 seconds.
No solve approached the 240-second limit.  These figures are reported only to
indicate scale; they are not assumptions in any certificate.

The reproduction uses CPython~3.13 because the audited binary modules are its
CPython~3.13 builds.  The pinned integer solver is OR-Tools 9.15.6755, loaded through
\file{computation/runtime_dependencies.py}.  SymPy is used only for the exact
ten-point polynomial calculations.  Floating-point feasibility, random
search, and a solver status other than explicit infeasibility are never used
to exclude a case.
